\section{符号说明}
表\ref{tab:notations}列出三问共用的主要符号。上标$t,a,v$分别表示文本、语音与视觉；原始素材时间以秒计，官方对齐序列长度以锚点数计。问题一采用自提取特征，问题二与问题三采用官方对齐特征，各问另行给出具体维度。问题三的可用模态集合$\mathcal A_i$、Shapley贡献$\phi_i^m$、作用份额$\omega_i^m$、主要参考模态$m_i^*$及局部删除效应$\Delta_i^m(E)$在第\ref{sec:q3}节定义。
\begin{table}[htbp]
\centering
\caption{主要符号及含义}
\label{tab:notations}
\begin{tabularx}{\linewidth}{p{3.5cm}X}
\toprule
符号 & 含义 \\
\midrule
$i,u,k,j$ & 样本、原词、输出位置及特征帧的索引 \\
$m\in\{t,a,v\}$ & 模态索引；涉及帧池化时$m\in\{a,v\}$ \\
$K,\ g_i(k)$ & 固定接口长度（$K=50$）及第$k$个内容词元的父词编号 \\
$P_{ik},\ O^m_{ik}$ & 内容位置掩码与模态观测掩码，取值均为0或1 \\
$A_{iu}=[a_{iu},b_{iu})$ & 第$i$条样本第$u$个原词的半开时间区间（s） \\
$J^m_{ij}=[s^m_{ij},e^m_{ij})$ & 模态$m$中第$j$个特征帧的时间支撑区间（s） \\
$\bm f^m_{ij},\ r^m_{ij}$ & 帧级特征向量及有效性标记 \\
$\ell^m_{ij}(A),\ \omega^m_{ij}(A)$ & 帧与锚点的有效交叠时长（s）及归一化池化权重 \\
$\bm x^m_{iu},\ X^m_i$ & 词级池化向量及固定50位的模态特征矩阵 \\
$c_{iu},\ d^m_{iu}$ & CTC路径置信度与词级模态覆盖指标，均位于$[0,1]$ \\
$B,\ s$ & 边界扰动次数与扰动标准差，分别为50次和0.04\,s \\
$V^m_{iu},\ \sigma_{im}^2$ & 扰动池化向量方差及样本内模态帧特征离散度 \\
$\widetilde V^m_{iu}$ & 由样本内尺度归一化的边界扰动方差 \\
$Q^m_{iu},\ \gamma_m$ & 综合质量分数及敏感度衰减系数，$\gamma_a=\gamma_v=1$ \\
$T_i^a,T_i^v,\Delta_i$ & 音轨、视频轨的解码跨度及其差$\Delta_i=T_i^a-T_i^v$（s） \\
$H^m_{ik},\widetilde O^m_{ik}$ & 问题二的人工遮挡掩码及遮挡后的观测掩码 \\
$\Omega_i,\rho^m_{ik}$ & 样本内观测候选集合及补全融合权重 \\
$y_i,r_i,\hat y_i,\hat r_i$ & 极性与强度真值及相应预测值 \\
$p_i^T,p_i^S,J$ & 教师、学生类别分布及验证集联合选模指标 \\
\bottomrule
\end{tabularx}
\end{table}
\FloatBarrier
